\documentclass[12pt]{article}
\usepackage[spanish]{babel}
\usepackage{lmodern,fontspec}
\setmainfont{Arial}
\usepackage{amsmath,amssymb,xcolor,booktabs,array,graphicx,pgfplots}
\pgfplotsset{compat=1.18}
\usepackage[paperwidth=841mm,paperheight=594mm,margin=22mm]{geometry}
\definecolor{ink}{HTML}{142E43}
\definecolor{teal}{HTML}{187B80}
\definecolor{muted}{HTML}{54636E}
\pagestyle{empty}
\setlength{\parindent}{0pt}
\setlength{\parskip}{4mm}
\setlength{\abovedisplayskip}{3mm}
\setlength{\belowdisplayskip}{3mm}
\DeclareMathSizes{22}{22}{16}{12}
\DeclareMathSizes{24}{24}{17}{13}
\DeclareMathSizes{34}{34}{24}{17}
\DeclareMathSizes{50}{50}{35}{25}
\newcommand{\head}[1]{{\color{teal}\fontsize{32}{38}\selectfont\bfseries #1\par}\vspace{1mm}}
\newcommand{\minor}[1]{{\fontsize{26}{32}\selectfont\bfseries #1\par}}
\newsavebox{\colA}\newsavebox{\colB}\newsavebox{\colC}
\begin{document}
\color{ink}
{\fontsize{58}{65}\selectfont\bfseries Una cota asintótica para Kuramoto\par}
{\fontsize{30}{37}\selectfont Siete clusters próximos a cuatro fases\hfill\textbf{Los Kuramotardos}\par}
{\fontsize{22}{28}\selectfont Investigathon 2026\hfill Optimización\par}
\vspace{2mm}
{\color{teal}\hrule height 1mm}
\vspace{3mm}
{\fontsize{50}{58}\selectfont\color{teal}
\[\boxed{\displaystyle\limsup_{\ell\to\infty}\mu(T_\ell)<\frac{3467}{5000}=0{,}6934}\]}
{\fontsize{24}{31}\selectfont\centering
Equilibrio exacto y estabilidad estricta en cada plantilla; siete masas límite positivas y fases límite $(0,0,\pi/2,\pi/2,\pi,\pi,3\pi/2)$.\par}
\vspace{7mm}
\fontsize{24}{31}\selectfont
\sbox{\colA}{\begin{minipage}[t]{250mm}
\head{1. Modelo y clase estudiada}
Partimos del Kuramoto homogéneo sobre grafos simples no dirigidos. Para estudiar grupos de fase usamos \textbf{plantillas ponderadas}: masas $x_i>0$, $\sum_i x_i=1$ y densidades $P=P^T\in[0,1]^{7\times7}$.
\[
F_i=\sum_j x_jP_{ij}\sin(\theta_i-\theta_j)=0,
\qquad \mu(T)=\min_i(Px)_i.
\]
La energía es $E=\sum_{i<j}x_ix_jP_{ij}[1-\cos(\theta_i-\theta_j)]$ y $H=\nabla^2E$. Se exige $H|_{\mathbf1^\perp}\succ0$ \textbf{después} de comprobar $F=0$.

\begin{center}
\begin{tikzpicture}
\begin{axis}[width=190mm,height=125mm,axis equal image,
 xmin=-1.6,xmax=1.6,ymin=-1.55,ymax=1.55,
 axis lines=middle,axis line style={muted!50},
 ticks=none,clip=false,
 xlabel={$\cos\theta$},ylabel={$\sin\theta$},
 xlabel style={at={(axis description cs:1,0.5)},anchor=west,font=\fontsize{20}{24}\selectfont},
 ylabel style={at={(axis description cs:0.5,1)},anchor=south,font=\fontsize{20}{24}\selectfont}]
\addplot[muted!55,thin,domain=0:360,samples=121] ({cos(x)},{sin(x)});
\addplot[only marks,mark=*,mark size=5pt,teal] coordinates {(1,0) (0,1) (-1,0) (0,-1)};
\node[anchor=south west,font=\fontsize{22}{27}\selectfont] at (axis cs:1.05,0.08) {$A$: 2};
\node[anchor=south west,font=\fontsize{22}{27}\selectfont] at (axis cs:0.1,1.08) {$B$: 2};
\node[anchor=south east,font=\fontsize{22}{27}\selectfont] at (axis cs:-1.05,0.08) {$C$: 2};
\node[anchor=north west,font=\fontsize{22}{27}\selectfont] at (axis cs:0.1,-1.08) {$W$: 1};
\end{axis}
\end{tikzpicture}
\end{center}
{\fontsize{21}{27}\selectfont Las etiquetas cuentan clusters en cada fase límite; el tamaño de los puntos no representa masas.}

\minor{Qué permite el teorema}
Mitades desiguales, falta de reflexión, internos incompletos, adyacencias límite cero y distintas escalas de separación. No se suponen derivadas de la sucesión.

Se excluyen masas límite nulas, otros repartos y geometrías alejadas de estas cuatro fases. \textbf{No es una cota universal para la conectividad crítica.}
\end{minipage}}
\sbox{\colB}{\begin{minipage}[t]{250mm}
\head{2. Reducción y desigualdad clave}
\textbf{Del equilibrio a cuatro cortes.} Si $h$ es la máxima desviación angular, la identidad exacta de equilibrio y la positividad de $H$ fuerzan enlaces opuestos $O(h^2)$ y cuatro cancelaciones principales.

Los flujos nulos se excluyen por grado antes de tomar centros ponderados. Completar internos se usa solo como relajación; no como operación que preserve equilibrio finito. Queda un vector variacional $z\ne0$.

\textbf{Conservar las masas vecinas.} Normalizamos el flujo común a uno. Para una hoja dividida, $f_0+f_1=1$, $\tau_i=f_i/m_i$ y una hoja vecina de masa $B$ implican $\tau_i\leq B$. Entonces
\begin{align*}
B-\tau_0\tau_1(m_0+m_1)
&=m_0\tau_0(B-\tau_1)\\
&\quad+m_1\tau_1(B-\tau_0)\geq0.
\end{align*}
Sean $a,b,c$ los productos $\tau_0\tau_1$ de A, B, C. Las masas de hoja se normalizan como $y_i=M_i/N$, con $N=\sum_i M_i$:
\[
\boxed{\begin{aligned}
a y_A&\leq y_B,\ y_W,\\
b y_B&\leq y_A,\ y_C,\\
c y_C&\leq y_B,\ y_W.
\end{aligned}}
\]
\textbf{Grados y capacidades a la misma escala.} Suponiendo grado límite al menos $u$, propagar $0<L_i\leq y_i\leq1-u$. Con $D=\max_i L_i(u-L_i)$, cada corte adyacente $pq$, con otras hojas $r,s$, exige
\[
\frac{4}{(1-L_r-L_s)^2}\leq N^2\leq\frac{2}{D}.
\]
Las covarianzas cumplen $|\kappa|\leq\sqrt{ab}$, $|\kappa|+\sqrt{ab}\leq1$ y análogamente para $\lambda,b,c$.

\textbf{Estabilidad incompatible.} Dos pruebas necesarias dan $z^TL_1^TL_1z\leq\|z\|^2$ y $z^TL_2^TL_2z\leq\|z\|^2$. Basta certificar, para algún $\alpha\in[0,1]$,
\[
\alpha L_1^TL_1+(1-\alpha)L_2^TL_2-I\succ0,
\]
o excluir antes la caja por masas y capacidades. Prueba completa y matrices: véase la nota citada.
\end{minipage}}
\sbox{\colC}{\begin{minipage}[t]{250mm}
\head{3. Un certificado del continuo}
Para $u=3467/5000$, las condiciones de grado confinan los tres productos a
\[
(a,b,c)\in\left[\frac{967}{1533},\frac{1533}{967}\right]^3.
\]
Un árbol de \textbf{15.511 nodos} cubre el cubo y sus bordes. Todas sus \textbf{7.756 hojas} tienen una exclusión racional estricta:
\begin{center}
\renewcommand{\arraystretch}{1.35}
\begin{tabular}{@{}p{191mm}r@{}}
\toprule
Justificación & Hojas\\
\midrule
Masas y capacidades & 3.742\\
Matriz de comparación & 2.362\\
Sylvester y Bernstein racional & 1.652\\
\midrule
\textbf{Total} & \textbf{7.756}\\
\bottomrule
\end{tabular}
\end{center}
Un verificador con Python estándar reconstruye dominio, partición y desigualdades. No importa el constructor ni usa autovalores flotantes para aceptar hojas. La reducción analítica es parte de la prueba.

\minor{Controles y reproducción}
El paquete pasó cinco órdenes en una copia aislada. Seis alteraciones del certificado fueron rechazadas. Los testigos finitos se comprueban con equilibrio exacto antes de estabilidad. \textbf{Constructor y verificadores proceden del mismo agente.}

\minor{Qué queda abierto}
La construcción conocida de O'Keeffe aporta $\mu_0\approx0{,}69290770637$. Para el supremo $s_*$ de valores límite de esta clase:
{\fontsize{34}{41}\selectfont\color{teal}
\[\mu_0\leq s_*\leq0{,}6934.\]}
El ancho es menor que $0{,}000493$. No se prueba $s_*=\mu_0$.

Un punto racional de la relajación satisface sus desigualdades a $0{,}6932$: no aporta una plantilla ni un nuevo récord. Alcanzar $\mu_0$ por esta vía requiere restricciones adicionales.

\end{minipage}}
\typeout{POSTER-COLUMN-A-HEIGHT=\the\dimexpr\ht\colA+\dp\colA\relax}
\typeout{POSTER-COLUMN-B-HEIGHT=\the\dimexpr\ht\colB+\dp\colB\relax}
\typeout{POSTER-COLUMN-C-HEIGHT=\the\dimexpr\ht\colC+\dp\colC\relax}
\noindent\usebox{\colA}\hfill\usebox{\colB}\hfill\usebox{\colC}\par
\vfill
{\color{teal}\hrule height 0.7mm}
\vspace{2mm}
{\fontsize{19}{24}\selectfont
\textbf{Grafos y normalización.} Para grafos con equilibrio constante en siete grupos de proporciones límite positivas, las fases límite del teorema y Hessiano completo estricto, agregar por promedios da $\delta/n\leq\mu(T)$ y el mismo límite superior. No se requieren bloques regulares. $(\delta+1)/n$ y $\delta/(n-1)$ solo comparten esta cota al crecer $n$; no se igualan a orden finito.
\par
\textbf{Fuentes y prueba.} Bullo, \textit{Lectures on Network Systems}, ed. 1.6, cap. 17. O'Keeffe, \textit{Beyond 11/16}, arXiv:2610.21271v1, IV.3: construcción y extremo inferior conocido. Desarrollo y certificado: \texttt{cota\_flujos\_06934.md}; paquete \texttt{kit-cota-flujos-06934-v1.zip}.
\par
\textbf{IA y revisión.} Codex (OpenAI) realizó derivaciones, código, cálculos, controles y redacción. Borrador pendiente de revisión matemática humana y comprobación de novedad. Preparado el 22/10/2026; sin entrega al evento.
}
\end{document}
